\chapter*{Preface}
Some of the most beautiful results in mathematics are exact values: $\sum\frac1{n^2}=\frac{\pi^2}6$, the product of
Wallis, the Gaussian integral, $\zeta(-1)=-\frac1{12}$. Others are approximations so good that they are almost as
striking, like Stirling's formula or the arithmetic--geometric mean, which doubles the number of correct digits of
$\pi$ at each step. This volume of \emph{Order and Chance} is about the art of calculating: finding the exact value of
a sum, a product or an integral, and, when no exact value exists, an approximation whose error one can bound.

It starts with finite sums, which have a calculus of their own, and generating functions, which turn sequences into
functions (Chapters~1--3). Chapters~4 and~5 compute integrals, with the tricks of the physicists and their
justification, and the Gamma function. Three chapters follow one idea, Fourier analysis, in three settings: the
circle, the finite group $\Z/n\Z$, where it becomes the fast Fourier transform, and the line, where it becomes Poisson
summation and decides the sign of a Gauss sum (Chapters~6--8). Complex analysis then turns real integrals and sums into
residues, and leads to the zeta function and its functional equation (Chapters~9 and~10). The volume ends with
asymptotics, the art of being precise about being approximate, and with Gauss's arithmetic--geometric mean
(Chapters~11 and~12).

The goal is a toolkit for the analytic number theory of Volume~5 and the probability of the later volumes, and the
pleasure of seeing very different methods agree on the same number.

Every theorem, proposition and lemma is proved; results quoted only for their interest appear as remarks, with a
reference to where they are proved. Exercises are graded \textbf{$\star$} for a warm-up, \textbf{$\star\star$} for a good
afternoon and \textbf{$\star\star\star$} for gems, and every one has a hint and a complete solution at the end of its
chapter, beginning with \emph{How to find it} before the \emph{Write-up}. ``Pause and think'' boxes ask you to guess
before you read on. The only fact used without proof is the Fubini--Tonelli theorem for general integrands; it is needed
rarely, since Chapter~4 proves Fubini's theorem for continuous functions on rectangles, and each use is labelled
\emph{Admitted (proved in Volume~6)}.

The series asks whether pure order and pure chance can be told apart, and why the primes look random yet obey laws. This floor is the scaffolding: series, integrals, and the functional equation of $\zeta$. The next floor takes up the primes in bulk, where a law and a Gaussian meet. The zeros are not yet carrying the primes.
